% ---- pg-362 （书页 350） ----

总方差的成分. 然而这个经验准则没有多少理论支持, 不应盲目运用. 如我们已经提到过的, 碎石图对选择合适的成分个数也是有用的.

\noindent\textbf{例 8.5}\quad (标准化数据的样本主成分)

5 支股票(JP Morgan, Citibank, Wells Fargo, Royal Dutch Shell 和 ExxonMobil)从 2004 年 1 月到 2005 年 12 月的周回报率数据取自纽约证券交易所. 周回报率定义为(当周收盘价 $-$ 上周收盘价)/上周收盘价, 并根据拆股和分红进行调整. 数据列于练习中的\hyperref[tab:8.4]{表 8.4}. 连续 103 周的观测显示是独立分布的, 但是股票之间的回报率是相关的, 因为如所预期的, 股票对一般经济条件的反应趋于一致.

用 $x_1,x_2,\cdots,x_5$ 分别记 5 支股票(JP Morgan, Citibank, Wells Fargo, Royal Dutch Shell, 和 ExxonMobil)的周回报率, 则
\[
\bar{\mathbf{x}}'=[0.001\ 1,0.000\ 7,0.001\ 6,0.004\ 0,0.004\ 0]
\]
且
\[
\mathbf{R}=\begin{bmatrix}
1.000&0.632&0.511&0.115&0.155\\
0.632&1.000&0.574&0.322&0.213\\
0.511&0.574&1.000&0.183&0.146\\
0.115&0.322&0.183&1.000&0.683\\
0.155&0.213&0.146&0.683&1.000
\end{bmatrix}
\]
注意 $\mathbf{R}$ 是标准化观测
\[
z_1=\frac{x_1-\bar{x}_1}{\sqrt{s_{11}}},\ z_2=\frac{x_2-\bar{x}_2}{\sqrt{s_{22}}},\ \cdots,\ z_5=\frac{x_5-\bar{x}_5}{\sqrt{s_{55}}}
\]
的协方差矩阵. 由计算机得到 $\mathbf{R}$ 的特征值和特征向量为
\[
\begin{aligned}
\hat{\lambda}_1&=2.437,&\qquad \hat{\mathbf{e}}_1'&=[\ 0.469,\ 0.532,\ 0.465,\ 0.387,\ 0.361]\\
\hat{\lambda}_2&=1.407,&\qquad \hat{\mathbf{e}}_2'&=[-0.368,-0.236,-0.315,\ 0.585,\ 0.606]\\
\hat{\lambda}_3&=0.501,&\qquad \hat{\mathbf{e}}_3'&=[-0.604,-0.136,\ 0.772,\ 0.093,-0.109]\\
\hat{\lambda}_4&=0.400,&\qquad \hat{\mathbf{e}}_4'&=[\ 0.363,-0.629,\ 0.289,-0.381,\ 0.493]\\
\hat{\lambda}_5&=0.255,&\qquad \hat{\mathbf{e}}_5'&=[\ 0.384,-0.496,\ 0.071,\ 0.595,-0.498]
\end{aligned}
\]
用标准化变量, 我们得到前两个样本主成分:
\[
\hat{y}_1=\hat{\mathbf{e}}_1'\mathbf{z}=0.469z_1+0.532z_2+0.465z_3+0.387z_4+0.361z_5
\]
\[
\hat{y}_2=\hat{\mathbf{e}}_2'\mathbf{z}=-0.368z_1-0.236z_2-0.315z_3+0.585z_4+0.606z_5
\]
这两个主成分的累计方差占总方差的百分比为
\[
\left(\frac{\hat{\lambda}_1+\hat{\lambda}_2}{p}\right)100\%=\left(\frac{2.437+1.407}{5}\right)\times 100\%=77\%
\]
这两个主成分有个有趣的解释. 第一主成分中 5 支股票有大致相同的加权, 或“指标”. 这个主成分可以解释为一般股票—市场成分, 或简称为市场成分.

第二主成分表示银行股(JP Morgan, Citibank, Wells Fargo)与石油股(Royal Dutch Shell, ExxonMobil)之间的对照(译者注: 在第二主成分中, 银行股的权有负号, 而石油股的权有正号, 且绝对值大致相同). 这个主成分可称为行业成分. 这样, 我们看到: 这些股票回报率的变
